\documentclass[11pt]{article}
\usepackage[a4paper,margin=1in]{geometry}
\usepackage{fontspec}
\usepackage[boldfont=false]{xeCJK}
\setCJKmainfont{FandolSong-Regular.otf}
\usepackage{amsmath}
\usepackage{amssymb}
\usepackage{array}
\usepackage{booktabs}
\usepackage{graphicx}
\usepackage{adjustbox}
\usepackage{enumitem}
\setlist[itemize]{topsep=2pt,partopsep=0pt,parsep=0pt,itemsep=2pt,leftmargin=1.4em}
\setlist[enumerate]{topsep=2pt,partopsep=0pt,parsep=0pt,itemsep=2pt,leftmargin=1.6em}
\usepackage{parskip}
\usepackage{fvextra}
\fvset{breaklines=true,breakanywhere=true,fontsize=\small}
\setlength{\parindent}{0pt}

\begin{document}

\begin{center}
  {\LARGE\bfseries States of matter}\\[0.35em]
  {\large Igcse Chemistry}
\end{center}
\vspace{0.6em}

\section*{The three states of matter}

Everything around you is made of tiny \textbf{particles} 粒子 — these can be \textbf{atoms} 原子, \textbf{molecules} 分子, or \textbf{ions} 离子. The \textbf{kinetic particle theory} 粒子动理论 says these particles are always moving. How close the particles are, how they are arranged, and how they move decides whether matter is a \textbf{solid} 固体, a \textbf{liquid} 液体, or a \textbf{gas} 气体.

\subsection*{Properties you can observe}

You do not need a microscope to tell the three states apart. They behave in different ways:

\begin{itemize}
\item A solid has a fixed shape and a fixed \textbf{volume} 体积. It does not flow and you cannot \textbf{compress} 压缩 it (squeeze it smaller).

\item A liquid has a fixed volume but no fixed shape. It flows and takes the shape of its container. It is almost impossible to compress.

\item A gas has no fixed shape and no fixed volume. It flows and spreads out to fill the whole container. A gas is easy to compress.

\end{itemize}

The table below sums up these properties. \textbf{Density} 密度 means how much \textbf{mass} 质量 is packed into a given volume.

\begin{center}
\adjustbox{max width=\linewidth}{%
\begin{tabular}{llll}
\toprule
\textbf{Property} & \textbf{Solid} & \textbf{Liquid} & \textbf{Gas} \\
\midrule
Shape & fixed & takes the shape of the container & fills the whole container \\
Volume & fixed & fixed & fills the whole container \\
Can it be compressed? & no & almost none & yes, easily \\
Does it flow? & no & yes & yes \\
Density & high & high & low \\
\bottomrule
\end{tabular}}
\end{center}

\subsection*{The particle picture}

The kinetic particle theory explains these properties by looking at three things: the \textbf{separation} 间距 of the particles (how far apart they are), their \textbf{arrangement} 排列 (the pattern), and their \textbf{motion} 运动 (how they move).

\begin{center}
\adjustbox{max width=\linewidth}{%
\begin{tabular}{llll}
\toprule
\textbf{} & \textbf{Solid} & \textbf{Liquid} & \textbf{Gas} \\
\midrule
Separation & touching, very close & close together & far apart \\
Arrangement & \textbf{regular} 规则 pattern & random, no pattern & random, no pattern \\
Motion & \textbf{vibrate} 振动 about fixed positions & move and slide past each other & move quickly in all directions \\
\bottomrule
\end{tabular}}
\end{center}

Strong \textbf{forces of attraction} 引力 hold the particles together. In a solid these forces are strong enough to hold every particle in place, so a solid keeps its shape. In a liquid the forces are weaker, so particles can move around. In a gas the particles move so fast that the forces hardly act at all, so the gas spreads out.

\section*{Changes of state}

When you heat or cool a substance, it can change from one state to another. You must know the name of each change.

\begin{center}
\adjustbox{max width=\linewidth}{%
\begin{tabular}{lll}
\toprule
\textbf{Change} & \textbf{What happens} & \textbf{Name} \\
\midrule
solid → liquid & \textbf{melting} 熔化 & melting \\
liquid → solid & \textbf{freezing} 凝固 & freezing \\
liquid → gas (at the surface, below the boiling point) & \textbf{evaporating} 蒸发 & evaporation \\
liquid → gas (all through the liquid) & \textbf{boiling} 沸腾 & boiling \\
gas → liquid & \textbf{condensing} 凝结 & condensation \\
\bottomrule
\end{tabular}}
\end{center}

A pure solid melts at one fixed temperature, the \textbf{melting point} 熔点. A pure liquid boils at one fixed temperature, the \textbf{boiling point} 沸点. The same substance freezes at its melting point and condenses at its boiling point.

\subsection*{Explaining changes of state with the particle theory}

Each change of state is really a change in the \textbf{energy} 能量 of the particles.

\begin{itemize}
\item When you heat a solid, the particles gain energy and \textbf{vibrate} faster. At the melting point the particles have enough energy to break away from their fixed places and slide around — the solid melts.

\item When you heat a liquid, the particles move faster. At the boiling point they have enough energy to fully escape the forces of attraction and become a gas.

\item Cooling does the opposite. The particles lose energy, move more slowly, and the forces of attraction pull them back together.

\end{itemize}

During a change of state the energy goes into breaking the forces of attraction, not into making the particles move faster. This is why the temperature stays the same while a substance is melting or boiling.

\subsection*{Heating and cooling curves}

A \textbf{heating curve} 加热曲线 is a graph of temperature against time as you heat a substance steadily. A \textbf{cooling curve} 冷却曲线 is the same graph as the substance cools.

On a heating curve there are two flat (level) parts:

\begin{itemize}
\item The first flat part is at the \textbf{melting point}. Here solid and liquid are both present. The heat energy breaks the forces holding the solid together, so the temperature does not rise.

\item The second flat part is at the \textbf{boiling point}. Here liquid and gas are both present, and the temperature again stays constant.

\end{itemize}

A cooling curve is the reverse. It has a flat part at the boiling point (the gas condenses) and a flat part at the melting point (the liquid freezes). As the particles slow down, \textbf{thermal energy} 热能 is released to the surroundings.

\section*{Gases: temperature, pressure and volume}

A gas pushes on the walls of its container. This push, spread over the area of the wall, is the \textbf{pressure} 压强 of the gas. Pressure comes from the gas particles hitting the walls.

\subsection*{Effect of temperature}

If you heat a fixed mass of gas while keeping the pressure the same, its volume increases.

Using the particle theory: heating gives the particles more \textbf{kinetic energy} 动能, so they move faster. They hit the walls harder and more often. To keep the pressure the same, the gas must take up more space, so the volume gets bigger.

If instead the volume is fixed (a sealed, rigid container), heating the gas makes the pressure rise, because the faster particles \textbf{collide} 碰撞 with the walls harder and more often.

\subsection*{Effect of pressure}

If you increase the pressure on a fixed mass of gas while keeping the temperature the same, its volume decreases. Squeezing the gas into a smaller space means the particles hit the walls more often, which is what a higher pressure means.

\section*{Diffusion}

\textbf{Diffusion} 扩散 is the spreading of particles from a region where they are crowded to a region where they are spread out — that is, from high \textbf{concentration} 浓度 to low concentration. It happens because particles are always moving in random directions.

Diffusion explains why you can smell food from across a room: the smell particles move and mix with the air until they reach your nose. Diffusion happens in gases and in liquids, but not in solids, because solid particles cannot move from place to place.

\subsection*{Rate of diffusion and molecular mass}

Lighter gas particles move faster than heavier ones at the same temperature. So a gas with a smaller \textbf{relative molecular mass} 相对分子质量 (a smaller mass for each molecule) has a faster \textbf{rate} 速率 of diffusion.

A classic experiment shows this. Cotton wool soaked in \textbf{ammonia} 氨气 ($\text{NH}_3$) is put at one end of a long glass tube. Cotton wool soaked in \textbf{hydrogen chloride} 氯化氢 ($\text{HCl}$) is put at the other end. Both gases diffuse along the tube and meet to form a white ring of \textbf{ammonium chloride} 氯化铵 ($\text{NH}_4\text{Cl}$).

\[
\text{NH}_3 + \text{HCl} \rightarrow \text{NH}_4\text{Cl}
\]

Ammonia has $M_r = 17$ and hydrogen chloride has $M_r = 36.5$. Ammonia is lighter, so it diffuses faster and travels further along the tube. The white ring forms nearer the hydrogen chloride end.


\end{document}
